\section{Bimanual Coordination, Handover, and Cooperative Tool Use}
\label{sec:bimanual}

\subsection{Bimanual learning is broader than tactile coordination}

Bimanual manipulation has grown rapidly through simulation benchmarks,
retargeting, imitation, and scalable demonstration generation. Bi-DexHands
provides a broad simulated reinforcement-learning benchmark
~\cite{chen2022bidexhands}. Bimanual Dexterity for Complex Tasks and related real
systems demonstrate increasingly capable coordinated policies
~\cite{shaw2024bidex}. DexMimicGen generates bimanual data from demonstrations
~\cite{jiang2024dexmimicgen}; ManipTrans transfers behaviours through residual
learning~\cite{li2025maniptrans}; object-centric human-motion approaches organize
demonstrations around manipulated objects~\cite{chen2024objectcentric}. These
works establish task and data scalability. They do not automatically test whether
a tactile latent explains cross-hand forces or contact transitions.

Late-August work further raises the vision/proprioception-only baseline.
Training-time modality masking improves robustness of a query-based bimanual VLA
without adding a tactile relation model~\cite{cheng2026modalitymasking};
PartialBiGrasp infers hidden local geometry from partial point clouds and searches
for force-closure-compatible grasp pairs~\cite{kaura2026partialbigrasp}. CLAP
learns a cross-embodiment action-conditioned video world model with bimanual data
among its supported settings~\cite{liu2026clap}, while rapid on-robot juggling
learns local dynamic models around a global prior under a reachable safe-set
constraint~\cite{lee2026robotjuggling}. Together they explain more bimanual
success through visual geometry, robust fusion, video prediction, and rapid
dynamics adaptation. A tactile representation must therefore demonstrate an
incremental benefit under occlusion, friction change, slip, or simultaneous load
transfer, not rely on bimanual task success alone.

This distinction matters for literature positioning. The broad claim ``bimanual
manipulation is underexplored'' is no longer defensible. The narrower intersection
of \emph{observation-grounded tactile structure}, \emph{explicit interface/load
targets}, and \emph{matched representation comparisons} remains much less
populated in the mapped corpus. A contribution should therefore be framed around
that intersection, not around the existence of two hands.

Bi-Touch is an important early direct baseline: it used a concatenated tactile
policy state and sim-to-real reinforcement learning for contact-rich bimanual
pushing, reorientation, and gathering tasks~\cite{bitouch2023}. It shifted the question from whether
touch can help bimanual control to whether an explicit relational representation
adds value beyond a matched tactile policy input. A new graph method should
therefore include an equivalent concatenation baseline even if the original
Bi-Touch implementation is not the primary reproduction target.

\subsection{Roles, symmetry, and mechanical necessity}

A bimanual taxonomy can distinguish symmetric and asymmetric actions, shared and
separate objects, coordinated and independent motions, and temporal role changes
~\cite{krebs2022taxonomy}. For representation research, two additional questions
are decisive:

\begin{enumerate}
  \item \textbf{Is coordination mechanically necessary?} If one hand can complete
  the task while the other supplies a decorative or easily ignored motion, task
  success does not diagnose cross-hand representation.
  \item \textbf{Can the roles be exchanged?} A representation that encodes fixed
  left/right identities may fail when actor and stabilizer swap, even though the
  physical relation is unchanged.
\end{enumerate}

BUDS makes stabilizer--actor coordination explicit and demonstrates that one hand
can create conditions for the other~\cite{grannen2023stabilize}. This is a useful
conceptual bridge from shared-object manipulation to tool--workpiece tasks.
However, role assignment should be coupled to measurable necessity: removing,
delaying, or perturbing the stabilizing hand ought to change forces, error, or
success in a predictable way.

\subsection{Handover as a load-transfer benchmark}

Handover provides an unusually clean phase structure. Human-inspired and
robot--robot/robot--human controllers have long used load, pull, grip force, or
slip cues to govern transfer and release~\cite{costanzo2021handover}. Contact-aware
dexterous handover further emphasizes when and where contacts occur
~\cite{wang2024contacthandover}. VTDexManip supplies a reproducible simulated
bimanual handover task with visual-tactile observations
~\cite{liu2025vtdexmanip}; recent tactile-reactive systems expand the scale and
realism of bimanual touch~\cite{niu2026trex}.

The most informative handover interval is dual contact. Let $F_L(t)$ and $F_R(t)$
be gravity-axis support estimates expressed in a common object frame. A simple
load-share target is

\begin{equation}
\rho_R(t)=\frac{\max(0,F_R(t))}
{\max(0,F_L(t))+\max(0,F_R(t))+\epsilon}.
\label{eq:loadshare}
\end{equation}

This ratio is interpretable and less scale-sensitive than absolute force, but it
is only valid if the force estimates are independently calibrated and frame
aligned. Contact phase, time-to-support, future receiver wrench, giver slip risk,
and release readiness provide complementary targets. The point is not to impose
one scalar definition; it is to make the transfer mechanism measurable.

Dynamic throw--catch handover explores timing, ballistic motion, and impact
recovery~\cite{huang2023dynamichandover}. It should be a later robustness task
rather than the primary cross-hand load benchmark, because the flight phase
removes the simultaneous object-mediated path. A quasi-static transfer with a
speed perturbation, object-mass change, or receiver delay offers cleaner causal
tests first.

Robot juggling reinforces this sequencing: impressive rapid online adaptation in
a dynamic bimanual task can be achieved from onboard vision and a local dynamics
model~\cite{lee2026robotjuggling}. It is therefore a strong adaptation reference,
but its flight/contact alternation does not provide the persistent, simultaneous
object-mediated load path needed to validate cross-hand tactile structure.

\subsection{Cooperative tool--workpiece manipulation}

One hand holding a brush while the other holds a plate is closely aligned with the
same representation problem, but its topology is more demanding. The actor hand
senses its grasp on the tool; the stabilizer hand senses the workpiece; the
working interface between brush and plate is remote from both local grasp
interfaces. Force and motion propagate through at least three contacts. This task
therefore tests whether a representation can move from
$H_L\!\leftrightarrow O\!\leftrightarrow H_R$ to
$H_L\!\leftrightarrow T\!\leftrightarrow W\!\leftrightarrow H_R$.

The adjacent literature supplies components rather than a saturated direct
benchmark. Tactile Tool Manipulation learns tool-related control with tactile
feedback~\cite{shirai2023tactiletool}; compliant tool--environment dynamics and
extrinsic contact-mode control model the working interface
~\cite{vandermerwe2022compliant,oller2024extrinsic}; RoboPack demonstrates
tactile-informed relational dynamics in tool-mediated dense contact
~\cite{ai2024robopack}; TACO covers tool--action--object semantics
~\cite{liu2024taco}; TAMEn targets tactile-aware closed-loop data collection in
contact-rich tasks such as dish washing~\cite{wu2026tamen}; and Semantic Contact
Fields make tool contacts explicit~\cite{ma2026semanticcontact}. These works make
the area adjacent, not empty. What remains scarce is a controlled study of
\emph{cross-hand tactile load propagation} with a held workpiece and matched
representation baselines. TAMEn is a direct system competitor for a dish-washing
demonstration; novelty cannot rest on performing that task with touch, but must
come from observation-grounded load-path inference, mechanism-level evidence, or
cross-topology transfer.

\begin{table}[t]
\centering
\caption{How the two principal task families complement each other. ``Mapped
evidence density'' is a qualitative inference from the curated corpus, not a
bibliometric claim.}
\label{tab:handover_tool}
\begin{tabularx}{\linewidth}{P{0.19\linewidth} Y Y}
\toprule
 & \textbf{Quasi-static handover} & \textbf{Cooperative tool--workpiece use} \\
\midrule
Core topology & hand--object--hand, changing support & hand--tool--workpiece--hand, persistent role asymmetry \\
Primary signal & load redistribution and release & working-contact wrench, stabilizer reaction, slip \\
Main advantage & clean phases and interpretable transition & harder multi-hop transfer and practical task relevance \\
Main confound & object pose/motion can reveal phase & visual trajectory, periodic motion, and compliance can reveal task state \\
Mapped evidence density & established handover/control literature; few matched representation studies & established tool/contact systems; few matched cross-hand load-flow studies \\
Best research role & first controlled representation benchmark & gated second topology after mechanism is established \\
\bottomrule
\end{tabularx}
\end{table}

\subsection{Designing a diagnostic tool task}

A diagnostic task should make the second hand necessary while remaining simple
enough to label. A rigid stylus or wiping puck against a plate is preferable to a
deformable brush for an initial experiment: the contact location, normal, and
applied path can be controlled, and simulation is more credible. The task can
then progress to a compliant brush once sensing and analysis are stable.

Useful experimental factors include tool stiffness, workpiece mass, friction,
path curvature, contact-force target, stabilizer grasp, actor/stabilizer role
swap, and brief perturbations of either hand. A \emph{coordination-necessity
check} should compare two hands, fixed workpiece, passive second hand, and removed
second-hand tactile input. If the actor succeeds with no measurable cost when the
stabilizer is removed or frozen, the task is not a strong cross-hand benchmark.

The tool task should enter a research programme after the simpler handover study
has established three points: inferred structure improves at least one
transition/load target over matched unstructured baselines; the gain is not due
to privileged state; and structure remains useful under object/grasp holds-out.
Before those conditions, tool use is valuable as a frozen task specification and
offline proxy, but a full hardware campaign risks combining unsolved sensing,
calibration, control, and representation problems.

\subsection{Cross-topology transfer as the stronger claim}

Training and testing separately on handover and tool use would produce two task
results. A stronger representation claim asks whether a model pretrained to infer
interfaces and loads on handover transfers to the longer tool topology with
fewer labels or demonstrations. The test should distinguish:

\begin{itemize}
  \item \textbf{encoder transfer:} freeze the local tactile encoder and retrain
  only the relation head;
  \item \textbf{relation transfer:} preserve edge functions while instantiating a
  new entity topology;
  \item \textbf{full fine-tuning:} initialize all weights and measure data
  efficiency;
  \item \textbf{zero-shot probing:} test whether phase/load variables are linearly
  recoverable without task fine-tuning.
\end{itemize}

If only full fine-tuning helps, the model may provide generic initialization rather
than compositional topology transfer. If relation modules transfer under role
swap and added tool/workpiece nodes, the evidence for structured factorization is
substantially stronger.
